\newpage
\section*{Реферат}

Отчет по научно-исследовательской работе изложен на ? страницах, содержит ? иллюстраций, ? таблицы, ? использованных источника, ? приложения.

Ключевые слова: распределенные вычислительные системы, планировщик задач, slurm, имитационное моделирование, нормальное распределение, логнормальное распределение, анализ датасета, кластер, ssh.

Объектом исследования являются методы диспетчеризации и планирования вычислительных ресурсов в распределенных кластерных системах под управлением СУПО SLURM на базе операционной системы Linux в среде виртуализации.

Целью работы является разработка программного имитатора пользовательских заявок на основе статистического анализа исторических данных суперкомпьютера (вариант: UID 50430, JobName \texttt{run\_bwa.sh}) и исследование динамики прохождения очередей задач в реальном планировщике SLURM с использованием механизмов масштабирования времени.

Результатом разработки стал комплекс программно-управляемых скриптов на языках Python и Bash, обеспечивающий фильтрацию исходного датасета \texttt{acct\_0921-0923}, вычисление статистических параметров человеческой ошибки ($\mu = 3.8621$, $\sigma = 0.8489$), генерацию случайного потока задач с усечением границ и их автоматическую отправку в пакетный диспетчер. На базе развернутого 5-узлового виртуального кластера проведено экспериментальное тестирование системы, в ходе которого средствами SLURM была собрана и проанализирована итоговая статистика соотношения «план – факт» для верификации математической модели.